\chapter{Master Problem Bank: Galvanic Cells, Electrodes \& Electrochemical Series}
\label{chap:qb_galvanic_cells}

\begin{tcolorbox}[enhanced,colback=subtleblue,colframe=primarynavy,arc=2mm,boxrule=1pt,
    title=\textbf{\large Part II Organization: Galvanic Cells, Reference Electrodes \& EC Series}]
This part compiles all questions and quantitative problems on electrode potentials, standard hydrogen electrode (SHE), calomel, silver-silver chloride, Daniell cell construction, salt bridge mechanics, and the electrochemical series from all 8 books. Identical and redundant variants from other books are co-located beneath each primary problem with full source citations.
\end{tcolorbox}

\section{Subtopic 2.1: Electrode Potential \& Electrical Double Layer}

\begin{problembox}
\textbf{Q.2.1} \hfill \textbf{[Primary Source: Neeraj Kumar Part 1 Ex-I Q1 / GRB  12.11]}
\label{q:qb-c2-01}

A metal rod is dipped in a solution of its own ions. Its single electrode potential is independent of:
\begin{enumerate}[label=(\alph*)]
    \item Temperature of the solution
    \item Concentration of the metal ions in solution
    \item Surface area of the metal rod exposed to the solution
    \item Nature of the metal and solvent
\end{enumerate}

\tcbline
\textbf{\color{accentamber}Cross-Book Redundant / Identical Variants:}
\begin{itemize}[leftmargin=*,itemsep=2pt]
    \item \textbf{[Variant v1: Pearson Ex-I Q1, p. 1]} The electrode potential of a metal does not depend on: (a) area of electrode (b) concentration of ions (c) temperature (d) pressure in case of gas electrodes.
    \item \textbf{[Variant v2: Narendra Avasthi Level-1 Q76, p. 12]} An intensive property among the following is: (a) electrode potential $E$ (b) resistance $R$ (c) Gibbs free energy change $\Delta G$ (d) heat capacity.
    \item \textbf{[Variant v3: Cengage Theory p. 3]} Explain the origin of the potential difference at a metal-solution interface using the concept of Helmholtz electrical double layer and Stern's compact/diffuse model.
\end{itemize}
\end{problembox}

\begin{problembox}
\textbf{Q.2.2} \hfill \textbf{[Primary Source: Atkins Topic 6C.1 / GRB  12.16]}
\label{q:qb-c2-02}

Why is the potential of the Standard Hydrogen Electrode (SHE) defined as zero at all temperatures, whereas the standard reduction potential of the saturated calomel electrode (SCE) has a non-zero temperature coefficient ($\partial E/\partial T \approx -0.65\,\text{mV\,K}^{-1}$)?

\tcbline
\textbf{\color{accentamber}Cross-Book Redundant / Identical Variants:}
\begin{itemize}[leftmargin=*,itemsep=2pt]
    \item \textbf{[Variant v1: Neeraj Kumar Part 1 Ex-I Q7, p. 1]} A standard hydrogen electrode has zero electrode potential because: (a) hydrogen is easiest to oxidize (b) this electrode potential is adopted as zero by international thermodynamic convention (c) hydrogen has one electron (d) hydrogen is the lightest element.
    \item \textbf{[Variant v2: Narendra Avasthi Level-1 Q82, p. 13]} The calomel electrode is reversible with respect to: (a) $\ce{Hg2^2+}$ (b) $\ce{Hg^2+}$ (c) $\ce{Cl-}$ (d) $\ce{H+}$.
    \item \textbf{[Variant v3: Pearson Ex-I Q15, p. 2]} What is the composition and standard potential of the saturated calomel electrode at $25^\circ\text{C}$? (a) $+0.2422\,\text{V}$ (b) $+0.2800\,\text{V}$ (c) $+0.3338\,\text{V}$ (d) $0.0000\,\text{V}$.
\end{itemize}
\end{problembox}

\section{Subtopic 2.2: Daniell Cell \& Salt Bridge Dynamics}

\begin{problembox}
\textbf{Q.2.3} \hfill \textbf{[Primary Source: GRB  12.13 Solved Ex 33 / NA Level-1 Q88]}
\label{q:qb-c2-03}

For the Daniell cell $\ce{Zn(s)} \mid \ce{ZnSO4(aq, } 1\,\text{M}\ce{)} \parallel \ce{CuSO4(aq, } 1\,\text{M}\ce{)} \mid \ce{Cu(s)}$:
\begin{enumerate}[label=(\alph*)]
    \item Write the anode half-reaction, cathode half-reaction, and overall cell reaction.
    \item What is the function of the salt bridge, and why does current stop flowing if the salt bridge is removed?
    \item Why cannot \ce{KCl} be used in a salt bridge if the cell contains $\ce{Ag+}$ or $\ce{Pb^2+}$ ions?
\end{enumerate}

\tcbline
\textbf{\color{accentamber}Cross-Book Redundant / Identical Variants:}
\begin{itemize}[leftmargin=*,itemsep=2pt]
    \item \textbf{[Variant v1: Neeraj Kumar Part 1 Ex-I Q22, p. 2]} A salt bridge contains \ce{KNO3} in agar-agar gel because: (a) \ce{K+} and \ce{NO3-} have almost equal ionic velocities ($t_+ \approx t_- \approx 0.5$) (b) agar-agar conducts current (c) it prevents mixing of solutions (d) all of these.
    \item \textbf{[Variant v2: Pearson Ex-I Q12, p. 2]} What happens when a Daniell cell operates for a long period? (a) mass of Zn rod increases (b) concentration of $\ce{Cu^2+}$ increases (c) mass of Cu rod increases and $[\ce{Zn^2+}]$ increases (d) cell potential increases.
    \item \textbf{[Variant v3: Cengage DPP 3.2 Q5, p. 3]} In which of the following electrochemical cells is a salt bridge NOT required? (a) Lead-acid storage battery (b) Daniell cell (c) Standard cadmium cell (d) Concentration cell without transference.
\end{itemize}
\end{problembox}

\section{Subtopic 2.3: Electrochemical Series \& Feasibility of Reactions}

\begin{problembox}
\textbf{Q.2.4} \hfill \textbf{[Primary Source: Neeraj Kumar Part 1 Ex-I Q3 / PRSN Ex-I Q8]}
\label{q:qb-c2-04}

The standard reduction potentials are given as:
$E^\circ(\ce{Mg^2+/Mg}) = -2.37\,\text{V}$,
$E^\circ(\ce{Al^3+/Al}) = -1.66\,\text{V}$,
$E^\circ(\ce{Zn^2+/Zn}) = -0.76\,\text{V}$,
$E^\circ(\ce{Fe^2+/Fe}) = -0.44\,\text{V}$,
$E^\circ(\ce{Cu^2+/Cu}) = +0.34\,\text{V}$,
$E^\circ(\ce{Ag+/Ag}) = +0.80\,\text{V}$.
Which of the following processes will occur spontaneously under standard conditions?
\begin{enumerate}[label=(\alph*)]
    \item Storing aqueous \ce{CuSO4} in a zinc container
    \item Stirring an \ce{Al(NO3)3} solution with a copper spoon
    \item Precipitation of silver from \ce{AgNO3} solution upon adding copper turnings
    \item Evolution of \ce{H2} gas upon adding copper to dilute \ce{HCl}
\end{enumerate}

\tcbline
\textbf{\color{accentamber}Cross-Book Redundant / Identical Variants:}
\begin{itemize}[leftmargin=*,itemsep=2pt]
    \item \textbf{[Variant v1: GRB Objective Set-1 Q45, p. 70]} Which metal will displace hydrogen from dilute \ce{H2SO4}? (a) Ag (b) Cu (c) Fe (d) Au.
    \item \textbf{[Variant v2: Narendra Avasthi Level-1 Q92, p. 14]} Given $E^\circ(\ce{Cr^3+/Cr}) = -0.74\,\text{V}$, $E^\circ(\ce{MnO4-/Mn^2+}) = +1.51\,\text{V}$, $E^\circ(\ce{Cr2O7^2-/Cr^3+}) = +1.33\,\text{V}$, $E^\circ(\ce{Cl/Cl-}) = +1.36\,\text{V}$. The strongest oxidizing agent is: (a) $\ce{MnO4-}$ (b) $\ce{Cl-}$ (c) $\ce{Cr^3+}$ (d) $\ce{Mn^2+}$.
    \item \textbf{[Variant v3: Cengage Theory Concept App 1 Q4, p. 14]} Can an aqueous solution of $1\,\text{M } \ce{FeSO4}$ be stored in a nickel vessel? Given $E^\circ(\ce{Ni^2+/Ni}) = -0.25\,\text{V}, E^\circ(\ce{Fe^2+/Fe}) = -0.44\,\text{V}$.
\end{itemize}
\end{problembox}
